Set $c_0=|\operatorname{GL}_4(\mathbb F_3)|=24\,261\,120$,
so $v_2(c_0)=9$, and define
\[
 H_M(n)=c_0\sum_{z\in X}
 \frac{P_M(z)\chi_{h_*}(\operatorname{Nrd}z)}{|R_z^\times|}
       \sum_{x\in V(\mathbb Q),\ Q_B(x)=n}f(z^{-1}xz). \tag{O.4}
\]
Left multiplication embeds each $R_z^\times$ into
$\operatorname{GL}_4(\mathbb Z)$; reduction modulo three is
injective on finite subgroups.  Indeed, a prime-order element
$1+3^aA$ in its kernel, with $a\geq1$ and $A$ nonzero modulo three,
has nonzero first term in its prime-th power modulo $3^{a+1}$,
or modulo $3^{a+2}$ when its order is three, a contradiction.
Thus $|R_z^\times|\mid c_0$.
The right-conjugation identity proves independence of the right
representative.  On the left, rational reduced norms are positive
and the global quadratic idele character is trivial on them.
The transformation law below shows that each summand is an actual
holomorphic theta series with the same level and character;
$c_0/|R_z^\times|$ is an integer.
Consequently the coefficients lie in the fixed character ring
$\mathcal O$.  Changing any lift of $P_M$ by $2^M$ changes every
coefficient by an element of $2^M\mathcal O$.  In particular
there is no loss of precision from an automorphism denominator.

